\chapter{Quorums \small{\textsf{DRAFT}}}\label{chapter.quorums}

\section*{Problems}
\begin{problems}
    \item Recall the longest chain PoS protocol, Ouroboros Praos, we studied in class.
        Ouroboros Praos is \emph{available} but not \emph{accountable}. We modify Ouroboros Praos to
        include a slashing mechanism: if the same key is used to create two blocks for
        the same height, we slash the key. Is the modified protocol $t$-accountably
        safe? If yes, prove it and find $t$. If not, is there another slashing mechanism
        that could make the protocol $t$-accountably safe?
    
    \item In this problem, we will consider the Simplex protocol run by $n$ permissioned
        nodes, with a tunable quorum size $q$ for notarization.
        \begin{enumerate}
            \item Suppose we want to maximize the resilience of the protocol, i.e.,
            maximize the number of adversarial nodes $f$ that can be tolerated such
            that the protocol is both safe and live. Derive the optimal quorum size
            $q$ and show that the resulting optimal resilience is $f = n/3$.
            \item Suppose you believe that an attack against safety is more likely than
            an attack against liveness (since a double-spend can provide significant
            rewards to the attacker). Hence, you want to tune Simplex to increase the
            resilience against safety attacks, even at the expense of decreasing the
            resilience against a liveness attack. Can this be done? If so, exhibit and
            plot the tradeoff between the two resiliences. If not, explain why not.
        \end{enumerate}
    
    \item Consider the Simplex protocol we presented in class. Consider the following
        variations of the protocol and in each case explain if the security (safety and
        liveness) of the protocol is preserved. In each case, provide a proof that the
        protocol retains each virtue or a show a counterexample in which it fails.
        \begin{enumerate}
            \item We completely remove the timer of the protocol.
            \item We change the timer from $3\Delta$ to $2\Delta$.
            \item We change the timer from $3\Delta$ to $4\Delta$.
        \end{enumerate}

    \item We modify the Simplex protocol to include a timestamp $r$ in each non-dummy
        block. When an honest party creates a block, they put their clock's time as the
        timestamp $r$. In a valid chain, consecutive blocks have non-decreasing
        timestamps and no timestamp is in the future. We assume that honest parties have
        sychronized clocks and the network is synchronous.
        We are interested in how much the timestamp of the finalized tip of an
        honest party can deviate from their clock's time. Compute an upper bound $v$ such
        that any honest party's tip will not deviate from their clock more than $v$
        except with negligible probability. Your bound does not need to be tight.
    
    \item Consider the Simplex protocol we studied in class.
        \begin{enumerate}
            \item Simplex is said to have finality whenever the number of honest
                nodes exceeds $2n/3$, where $n$ is the total number of nodes. Explain
                what that means.
            \item Simplex is said to be {\em $1/3$-accountable}. Explain what
                that means.
            \item Now suppose a synchronous network with $\Delta = 1$ second and
                $n = 100$ nodes and the nodes are all honest.
                \begin{enumerate}
                    \item Compute the expected growth rate of the notarized
                        chain (in blocks per second).
                    \item Compute the expected growth rate of the finalized
                        chain (in blocks per second).
                \end{enumerate}
            \item Now suppose $20$ nodes are still adversary, but instead of having the
                full $80$ honest nodes online, $30$ of them decide not to participate in
                the protocol and went on vacation to Ikaria. However, the protocol
                designer does not know this and the protocol parameters are not
                adjusted.
                \begin{enumerate}
                \item Compute a tight lower bound on the expected growth
                        rate of the longest notarized chain.
                \item Compute a lower bound on the expected growth rate of
                        the finalized chain. (Your lower bound does not need to be tight.
                        However if the lower bound is not tight, it must be positive.)
                \item Is the protocol safe? Is the protocol live?
                \end{enumerate}
            \item A quorum-based protocol is said to be {\em available} if it
                is safe \emph{and} live whenever the number of honest nodes online is
                greater than twice the number of adversary nodes. Is Simplex available?
        \end{enumerate}

\end{problems}

\section*{Further Reading}

{\color{red}
\begin{itemize}
\item Slashing
\end{itemize}
}
